\documentclass[../../main.tex]{subfiles}

\graphicspath{{../../image/}{../../figure/}}

\begin{document}

\section{习题}\label{section:ch08-习题}

\begin{exercise}\label{exercise:ch08-01}
（★）利用\eqref{align:ch08-预激活}、\eqref{align:ch08-激活}、\eqref{align:ch08-误差}，以及在\cref{theorem:ch08-反向传播方程}的证明中出现的链式分解，验证计算误差函数导数的反向传播公式［\eqref{align:ch08-隐藏误差}］。
\end{exercise}

\begin{exercise}\label{exercise:ch08-02}
（★★）采用一个多层网络，并从前向传播公式［\eqref{align:ch06-矩阵形式}］开始，用矩阵符号重写反向传播公式［\eqref{align:ch08-隐藏误差}］。请注意，结果涉及转置矩阵。
\end{exercise}

\begin{exercise}\label{exercise:ch08-03}
（★）使用Taylor展开，验证 $\mathcal{O}(\varepsilon)$ 在\eqref{align:ch08-中心差分}的右侧被抵消了。
\end{exercise}

\begin{exercise}\label{exercise:ch08-04}
（★★）考虑\cref{figure:ch06-两层网络}所示的两层网络，其中额外增加了与直接从输入到输出的跳层连接相对应的参数。通过拓展\cref{subsection:ch08-简单示例}所讨论的内容，写出误差函数关于这些额外参数的导数方程。
\end{exercise}

\begin{exercise}\label{exercise:ch08-05}
（★★★）在\cref{subsection:ch08-Jacobi矩阵}中，我们推导出了一种使用反向传播程序计算神经网络Jacobi矩阵的方法。请推导出一种基于前向传播方程计算Jacobi矩阵的替代形式。
\end{exercise}

\begin{exercise}\label{exercise:ch08-06}
（★★★）考虑双层神经网络并定义以下数据：
\begin{align}
\delta_k=\frac{\partial E_n}{\partial a_k},\qquad
M_{kk'}=\frac{\partial^{2}E_n}{\partial a_k\partial a_{k'}}.
\label{align:ch08-习题deltaM}
\end{align}
对于（i）两个权重都在第二层、（ii）两个权重都在第一层、（iii）每一层都有一个权重的元素，推导出用 $\delta_k$ 和 $M_{kk'}$ 表示的黑塞矩阵的元素表达式。
\end{exercise}

\begin{exercise}\label{exercise:ch08-07}
（★★★）将\cref{exercise:ch08-06}中关于双层神经网络的精确黑塞矩阵的结果进行扩展，使其包括直接从输入层到输出层的跳层连接。
\end{exercise}

\begin{exercise}\label{exercise:ch08-08}
（★★）使用平方和误差函数的神经网络黑塞矩阵的外积近似值由\eqref{align:ch08-外积近似}给出。将这一结果扩展到多个输出。
\end{exercise}

\begin{exercise}\label{exercise:ch08-09}
（★★）考虑如下形式的平方损失函数：
\begin{align}
E(\mathbf{w})=\frac{1}{2}\iint\left\{y(\mathbf{x},\mathbf{w})-t\right\}^{2}p(\mathbf{x},t)\,\mathrm{d}\mathbf{x}\,\mathrm{d}t,
\label{align:ch08-习题平方损失}
\end{align}
其中 $y(\mathbf{x},\mathbf{w})$ 为参数化函数，如神经网络。\cref{proposition:ch04-最优决策规则}表明，能使该误差最小化的函数 $y(\mathbf{x},\mathbf{w})$ 是由给定 $\mathbf{x}$ 情况下 $t$ 的条件期望给出的。该结果可证明 $E$ 相对于向量 $\mathbf{w}$ 的两个元素 $w_r$ 和 $w_s$ 的二阶导数为
\begin{align}
\frac{\partial^{2}E}{\partial w_r\partial w_s}=\int\frac{\partial y}{\partial w_r}\frac{\partial y}{\partial w_s}p(\mathbf{x})\,\mathrm{d}\mathbf{x}.
\label{align:ch08-习题平方损失二阶}
\end{align}
注意，对于 $p(\mathbf{x})$ 的有限样本，我们可以得到\eqref{align:ch08-外积近似}。
\end{exercise}

\begin{exercise}\label{exercise:ch08-10}
（★★）推导出具有单输出、逻辑斯谛 sigmoid 输出单元激活函数和交叉熵误差函数的网络的黑塞矩阵的外积近似表达式［\eqref{align:ch08-外积交叉熵}］，对应于采用平方误差函数时的结果［\eqref{align:ch08-外积近似}］。
\end{exercise}

\begin{exercise}\label{exercise:ch08-11}
（★★）推导出一个具有 $K$ 个输出、softmax 输出单元激活函数和交叉熵误差函数的网络的黑塞矩阵的外积近似表达式，对应于采用平方和误差函数时的结果。
\end{exercise}

\begin{exercise}\label{exercise:ch08-12}
（★★★）考虑矩阵等式
\begin{align}
\left(\mathbf{M}+\mathbf{v}\mathbf{v}^{\mathsf T}\right)^{-1}
=\mathbf{M}^{-1}-\frac{\left(\mathbf{M}^{-1}\mathbf{v}\right)\left(\mathbf{v}^{\mathsf T}\mathbf{M}^{-1}\right)}{1+\mathbf{v}^{\mathsf T}\mathbf{M}^{-1}\mathbf{v}}.
\label{align:ch08-伍德伯里}
\end{align}
它只是 Woodbury 恒等式［\eqref{align:appA-woodbury}］的一个特例。将\eqref{align:ch08-伍德伯里}应用于黑塞矩阵的外积近似表达式［\eqref{align:ch08-外积近似}］，推导出一个公式，该公式允许通过训练数据进行一次传播并用每个数据点更新黑塞矩阵，从而计算出黑塞矩阵的逆。需要注意的是，该算法可以采用 $\mathbf{H}=\alpha\mathbf{I}$ 进行初始化，其中 $\alpha$ 是一个小的常数，而且结果对 $\alpha$ 的精确值并不特别敏感。
\end{exercise}

\begin{exercise}\label{exercise:ch08-13}
（★★）验证\eqref{align:ch08-复合函数}的导数由\eqref{align:ch08-复合导数}给出。
\end{exercise}

\begin{exercise}\label{exercise:ch08-14}
（★★）逻辑斯蒂映射是由 $L_{n+1}(x)=4L_n(x)\left(1-L_n(x)\right)$ 和 $L_1(x)=x$ 的迭代关系定义的函数。写出 $L_2(x)$、$L_3(x)$、$L_4(x)$ 的求值轨迹方程，然后写出相应的导数 $L_1'(x)$、$L_2'(x)$、$L_3'(x)$、$L_4'(x)$ 的表达式。请勿简化这些表达式，但需要注意导数公式的复杂度是如何比函数本身的表达式的复杂度增长得更快的。
\end{exercise}

\begin{exercise}\label{exercise:ch08-15}
（★★）从示例函数\eqref{align:ch08-示例函数}的求值轨迹方程\eqref{align:ch08-求值轨迹}～\eqref{align:ch08-求值轨迹末}开始，利用\eqref{align:ch08-正切一般}推导出前向模式正切变量求值方程\eqref{align:ch08-正切轨迹}～\eqref{align:ch08-正切轨迹末}。
\end{exercise}

\begin{exercise}\label{exercise:ch08-16}
（★★）从示例函数\eqref{align:ch08-示例函数}的求值轨迹方程\eqref{align:ch08-求值轨迹}～\eqref{align:ch08-求值轨迹末}开始，并参考\cref{figure:ch08-求值轨迹}，利用\eqref{align:ch08-伴随链式}推导出逆模式伴随变量求值方程\eqref{align:ch08-伴随轨迹}～\eqref{align:ch08-伴随轨迹末}。
\end{exercise}

\begin{exercise}\label{exercise:ch08-17}
（★★★）考虑示例函数\eqref{align:ch08-示例函数}。写出 $\partial f/\partial x_1$ 的表达式，并在 $x_1=1$ 和 $x_2=2$ 时对该函数进行求值。使用求值轨迹方程\eqref{align:ch08-求值轨迹}～\eqref{align:ch08-求值轨迹末}求出变量 $v_1\sim v_7$，然后使用前向模式自动微分的求值轨迹方程求出正切变量 $\dot{v}_1\sim\dot{v}_7$，并确认求得的 $\partial f/\partial x_1$ 值与直接求得的 $\partial f/\partial x_1$ 值一致。同样，使用逆模式自动微分的求值轨迹方程求出伴随变量 $\bar{v}_1\sim\bar{v}_7$，并再次确认求得的 $\partial f/\partial x_1$ 值与直接求得的 $\partial f/\partial x_1$ 值一致。
\end{exercise}

\begin{exercise}\label{exercise:ch08-18}
（★★）将任意向量 $\mathbf{r}=(r_1,\cdots,r_D)^{\mathsf T}$ 表示为单位向量 $\mathbf{e}_i$ 的线性组合，其中 $i=1,\cdots,D$。证明可以通过设置 $\dot{\mathbf{x}}=\mathbf{r}$，只需要通过单次前向模式自动微分传播就可以求出形式为\eqref{align:ch08-Jacobi向量积}的函数与 $\mathbf{r}$ 的Jacobi矩阵的乘积。
\end{exercise}

\end{document}
